\begin{qfdefinition}[Agricultural markets for macro: corn, wheat, soybeans, cotton and coffee, crop calendars and sensitive stages, CBOT and ICE contracts, ethanol demand, and weather with the livestock cascade]\label{def:agricultural-markets}
\qffield{Statement} Gliner: agricultural supply depends on weather and planting, and ``some of the best agriculture traders have full-time meteorologists on staff''; crop calendars show the sensitive periods \citep[chap.~10, sec.\ ``Agriculture'', PDF p.~330]{s7-gliner}. \emph{Corn} (pp.~330--335): the United States produces over 35 percent and consumes over 30 percent of the world's corn and is the largest exporter; Iowa, Illinois and Nebraska produce 47 percent of U.S.\ corn; uses are feed, ethanol (27 percent of U.S.\ corn demand in 2011, after the Energy Policy Act of 2005) and food; Chinese feed imports rose from nothing in 2007 to over 5 percent of world imports; silking from June is the critical period; CBOT corn (C) is quoted in cents a bushel, 5,000 bushels a contract, 56 pounds a bushel, for March, May, July, September and December.

\emph{Cotton} (pp.~341--344): character, grade and staple; China, India and the United States 65 percent of production in 2013, China the largest producer and consumer, the United States the largest exporter with almost 30 percent; ICE cotton (CT) in cents a pound, 50,000 pounds, roughly 104 bales of 480 pounds. \emph{Coffee} (pp.~344--348): arabica about 60 percent of production, robusta the rest; Brazil the largest producer (37.2 percent in 2012--13) and Vietnam second (16.6 percent), in bags of 60 kilograms (Table~10.21); the United States and the European Union almost half of demand; ICE arabica (KC). \emph{Weather and livestock} (p.~333): when grain prices rise sharply, feeding livestock becomes expensive and farmers slaughter early, so near-term livestock prices fall while later prices rise.

\qffield{Notation} Futures month codes F (Jan), H (Mar), K (May), N (Jul), Q (Aug), U (Sep), V (Oct), X (Nov), Z (Dec); one bushel of corn 56 pounds, of wheat and soybeans 60.

\qffield{Conditions} Shares, contract terms and delivery months are Gliner's for 2012--2013 (USDA tables as he reproduces them); CBOT grain terms are compared with the December 2024 filing below. A corn contract at 500 cents is \$25,000 notional (checks, section~14).

\qffield{Practice} Translate a view on a crop's balance sheet (\hyperref[def:contango-backwardation-stocks]{ending stocks}) into a specific contract month, grade and delivery location; the basis between local cash and futures then matters.

\qffield{Pitfalls} \emph{Beyond these sources}: El Ni\~no and La Ni\~na, NOAA drought monitors, relatively inelastic demand, Russian wheat export bans, coffee's biennial cycle and 20-year tree life are plausible influences that these sources do not cover. Two against ICE: Coffee ``C'' (KC) is 37,500 pounds, quoted in cents a pound, for March, May, July, September and December, physically delivered; Cotton No.~2 (CT) is 50,000 pounds net weight, for March, May, July, October and December, matching H, K, N, V, Z \citep[contract specifications]{ice-softs-specs}.

\emph{CBOT grain terms}. Each contract is 5{,}000 bushels with a minimum fluctuation of a quarter cent a bushel (\$12.50 a contract). Corn: No.~2 yellow at par, No.~1 at 1\textonehalf\ cents over and No.~3 at 2 to 4 cents under, in September, December, March, May and July. Soybeans: No.~2 yellow at par, No.~1 at 6 cents over and No.~3 at 6 cents under, in September, November, January, March, May, July and August. Wheat: No.~2 Soft Red Winter, Hard Red Winter, Dark Northern Spring and Northern Spring at par and No.~1 at 3 cents over, in July, September, December, March and May.

All three stop trading on the business day before the 15th calendar day of the delivery month, and daily price limits are reset twice a year at 7 percent of a 45-day average settlement, rounded to 5 cents with a 20-cent floor for corn, expanding by 1.5 times after a limit settlement \citep[Exhibit, Rules 10101--10104, 11101--11102 and 14101--14102]{cbot-24-443}. The filing itself raised load-out premiums on shipping certificates from 6 to 9 cents a bushel and the St.\ Louis--Alton corn and soybean differential from 16.25 to 24 cents over par, effective after the 2027 expiries.

The current Chapters 10, 11 and~14 (official PDFs, undated) agree with the reprint on these terms, and carry the location differentials in two schedules, for corn up to December 2027 and from March 2028, and for soybeans up to November 2027 and from January 2028 \citep[Rules 10101--10102, 10105, 11101--11102, 11105 and 14101--14102]{cbot-grain-chapters-2026}; no change after the January 2025 effective date was found in the terms above. Gliner's months for corn and soybeans match; his cotton contract size ``50,00 0 lbs'' is a typo for 50,000.

\qffield{See also} \hyperref[def:contango-backwardation-stocks]{contango and ending stocks}; \hyperref[def:futures-hedge-basis]{futures hedges and basis}; \hyperref[def:commodity-indices-risk-return]{commodity indices}.

\qffield{Further reading} WASDE \citep{physical-usda-wasde-2026}.

\end{qfdefinition}
